\documentclass[11pt]{article}
\usepackage{fontspec}
\usepackage[french]{babel}
\usepackage{amsmath}
\usepackage{amssymb}
\usepackage{amsthm}
\usepackage{mathrsfs}
\usepackage{tikz-cd}
\usepackage{enumitem}
\usetikzlibrary{babel}

\newcommand{\printedpage}[1]{\marginpar{\footnotesize p.~#1}}
\newcommand{\msp}[1]{\marginpar{\footnotesize [#1]}}
\newcommand{\sX}{\mathscr{X}}
\newcommand{\tX}{\tilde{\mathscr{X}}}

\begin{document}

\printedpage{88}
% head: A. Grothendieck : Dérivateurs, Ch. XII, Édition des Universités de Montpellier II et Paris VII
\[
Q \quad = \qquad
\begin{tikzcd}
 & X_0 \arrow[dl, "\alpha_1"'] \arrow[dr, "\alpha_2"] & \\
X_1 \arrow[dr, "f_1"'] & & X_2 \arrow[dl, "f_2"] \\
 & Y, &
\end{tikzcd}
\]
on peut le compléter par
\[
f_0 \overset{\text{déf}}{=} f_1\alpha_1 = f_2\alpha_2
\]
en un diagramme
\[
\begin{tikzcd}[row sep=large, column sep=large]
 & X_0 \arrow[dl, "\alpha_1"'] \arrow[dr, "\alpha_2"] \arrow[dd, "f_0", ""'{name=F0, pos=0.65}] & \\
X_1 \arrow[dr, "f_1"', ""{name=F1, pos=0.4}] & & X_2 \arrow[dl, "f_2", ""'{name=F2, pos=0.4}] \\
 & Y, &
\arrow[from=F1, to=F0, bend left, "\mathrm{id}", shorten <=2pt, shorten >=2pt]
\arrow[from=F2, to=F0, bend right, "\mathrm{id}"', shorten <=2pt, shorten >=2pt]
\end{tikzcd}
\]
ce qui montre qu'on a une donnée du type $(*)$ ci-dessus, définissant donc une flèche
\[
\varphi_Q : \quad \sX = \int
\begin{tikzcd}[row sep=small, column sep=small]
 & X_1 \\
X_0 \arrow[ur] \arrow[dr] & \\
 & X_2
\end{tikzcd}
\qquad \longrightarrow \qquad Y
\]
qui donne lieu à un diagramme cubique

\msp{page 131}
\[
\begin{tikzcd}[row sep=small, column sep=small]
 & X_1 \arrow[dd, "\mathrm{id}"'] \arrow[rrd, "i_1", ""'{name=I1}] & & \\
X_0 \arrow[ru, "\alpha_1"] \arrow[rrr, dotted, "i_0", ""'{name=I0, pos=0.8}] \arrow[rrd, "\alpha_2"] \arrow[dd, "\mathrm{id}"'] & & & \sX \arrow[dd, "\varphi"] \\
 & X_1 \arrow[rrd, "f_1"] & X_2 \arrow[ru, "i_2"', ""{name=I2}] \arrow[dd, "\mathrm{id}"'] & \\
X_0 \arrow[ru, "\alpha_1"] \arrow[rrr, dotted, "f_0"] \arrow[rrd, "\alpha_2"] & & & Y, \\
 & & X_2 \arrow[ru, "f_2"'] &
\arrow[from=I1, to=I0, "\lambda_1"', shorten <=1pt, shorten >=1pt]
\arrow[from=I2, to=I0, bend left, "\lambda_2", shorten <=1pt, shorten >=1pt] % Resolved from PDF leaf 11: a solid curve, unlike the genuinely dotted i_0 line.
\end{tikzcd}
\]
% foot: Version 7 décembre 2002

\printedpage{89}
% head: A. Grothendieck : Dérivateurs, Ch. XII, Édition des Universités de Montpellier II et Paris VII
commutatif sauf le carré supérieur, i.e. on a
\[
\varphi i_1 = f_1, \qquad \varphi i_2 = f_2, \qquad \text{et de plus} \quad \varphi i_0 = f_0,
\]
$\varphi$ transformant de plus les homomorphismes fonctoriels
\[
\lambda_1 : i_1\alpha_1 \longrightarrow i_0, \qquad \lambda_2 : i_2\alpha_2 \longrightarrow i_0
\]
en les homomorphismes fonctoriels identiques pertinents.

\setcounter{footnote}{85}
\textbf{Définition 3.} On dit que le diagramme commutatif
\[
\begin{tikzcd}[row sep=small, column sep=small]
 & X_0 \arrow[dl, "\alpha_1"'] \arrow[dr, "\alpha_2"] & \\
X_1 \arrow[dr, "f_1"] & & X_2 \arrow[dl, "f_2"'] \\
 & Y &
\end{tikzcd}
\tag*{$(*)$}
\]
est \textit{$W$-homotopiquement cocartésien} (\footnote{ou simplement $W$-cocartésien, ou qu'il définit $Y$ comme somme amalgamée $W$-homotopique de etc.}) si l'homomorphisme canonique
\[
\sX = \int
\begin{tikzcd}[row sep=small, column sep=small]
 & X_1 \\
X_0 \arrow[ur] \arrow[dr] & \\
 & X_2
\end{tikzcd}
\qquad \xrightarrow{\ \varphi\ } \qquad Y
\]
qu'il définit est dans $W$. (\footnote{Introduire $\varphi$ dans une proposition en forme, et donner ses propriétés essentielles : flèches strictement cartésiennes sont transformées en flèches \textit{identiques} $+$ lemmes 1 et 2 des p. 152, 154.})

\msp{page 132}
On a alors le :

\textbf{Corollaire 1.} \textit{Supposons} W(7 bis)\textit{. Si le carré} $(*)$ \textit{est $W$-homotopiquement cocartésien, alors}
\[
\left\{
\begin{aligned}
\alpha_1 \in W &\quad\iff\quad f_2 \in W \\
\alpha_2 \in W &\quad\iff\quad f_1 \in W.
\end{aligned}
\right.
\]

Cela résulte aussitôt du diagramme cubique de la page précédente, où toutes les flèches verticales sont dans $W$, et de l'énoncé similaire concernant la face supérieure du cube, donné dans la proposition 13, p. 125.

\textbf{Corollaire 2.} \textit{Supposons} W(7)\textit{. Considérons un diagramme cubique}
% foot: Version 7 décembre 2002

\printedpage{90}
% head: A. Grothendieck : Dérivateurs, Ch. XII, Édition des Universités de Montpellier II et Paris VII
\[
\begin{tikzcd}[row sep=small, column sep=small]
 & X_1 \arrow[dd, "p_1"'] \arrow[rrd, "f_1"] & & \\
X_0 \arrow[ru, "\alpha_1"] \arrow[rrd, "\alpha_2"] \arrow[dd, "p_0"'] & & & Y \arrow[dd, "p"] \\
 & X'_1 \arrow[rrd, "f'_1"] & X_2 \arrow[ru, "f_2"] \arrow[dd, "p_2"'] & \\
X'_0 \arrow[ru, "\alpha'_1"] \arrow[rrd, "\alpha'_2"] & & & Y', \\
 & & X'_2 \arrow[ru, "f'_2"] &
\end{tikzcd}
\]
\textit{où les deux faces inférieure et supérieure du cube sont des carrés $W$-homotopiquement cocartésiens. Considérons les conditions}

\msp{page 133}
\begin{enumerate}[label={}]
\item[\textit{a)}] $p_0 \in W$\textit{.}
\item[\textit{b)}] $p_1$ \textit{et} $p_2 \in W$\textit{.}
\item[\textit{c)}] $p \in W$\textit{.}
\end{enumerate}

\textit{Alors deux de ces conditions impliquent la troisième.}

Résulte aussitôt de la définition 3, et du résultat similaire (cor. page 126) pour le diagramme cubique similaire, avec $Y$, $Y'$ remplacées par $\sX = \int
\begin{tikzcd}[row sep=small, column sep=small]
 & X_1 \\
X_0 \arrow[ur] \arrow[dr] & \\
 & X_2
\end{tikzcd}$
et $\sX' = \int
\begin{tikzcd}[row sep=small, column sep=small]
 & X'_1 \\
X'_0 \arrow[ur] \arrow[dr] & \\
 & X'_2
\end{tikzcd}$\;.

\section*{3.3\quad Intégration propre versus intégration lisse. Énoncé d'équivalence entre les deux points de vue. Commentaires sur le cas des limites $W$-homotopiques quelconques.}

La définition 3, et les corollaires précédents, n'ont d'intérêt que dans la mesure où on dispose de critères maniables, permettant d'affirmer qu'un carré commutatif est bel et bien homotopiquement cocartésien. C'est ce qu'on va essayer de faire tantôt. Il y a aussi une perplexité, concernant le rôle privilégié joué, dans la théorie développée ici, par les immersions fermées, à l'exclusion des immersions ouvertes. Nous avons en somme ramené, du
% foot: Version 7 décembre 2002
\printedpage{91}
% head: A. Grothendieck : Dérivateurs, Ch. XII, Édition des Universités de Montpellier II et Paris VII
\msp{page 134}
point de vue homotopique, les diagrammes quelconques $\begin{tikzcd}[row sep=small, column sep=small]
 & X_1 \\
X_0 \arrow[ur] \arrow[dr] & \\
 & X_2
\end{tikzcd}$ dans Cat, au cas d'un diagramme d'immersions fermées, où la somme amalgamée prend des allures familières. Mais pourquoi pas aussi bien le ramener au cas d'immersions ouvertes? On considérerait alors le diagramme de départ comme définissant une catégorie \textit{cofibrée} sur
\[
\Phi^{\circ} = \Psi = \
\begin{tikzcd}[row sep=small, column sep=small]
 & 1 \\
0 \arrow[ur] \arrow[dr] & \\
 & 2,
\end{tikzcd}
\]
soit $\tX$. On a dans $\Psi$ les deux ouverts
\[
\underbrace{\Psi/1}_{=\{0,1\}} = U_1, \qquad \underbrace{\Psi/2}_{=\{0,2\}} = U_2, \qquad U_1 \cap U_2 = \{0\},
\]
d'où dans $\tX$ deux ouverts images inverses
\[
\tX/1, \qquad \tX/2, \qquad \text{avec } \tX/1 \cap \tX/2 = X_0,
\]
d'où un diagramme commutatif d'inclusions ouvertes
\[
\begin{tikzcd}
 & X_0 \arrow[dl, "\alpha'_1"'] \arrow[dr, "\alpha'_2"] & \\
\tX/1 \arrow[dr, "i'_1"'] & & \tX/2 \arrow[dl, "i'_2"] \\
 & \tX, &
\end{tikzcd}
\tag*{$(*)$}
\]

\msp{page 135}
lequel reçoit le diagramme non commutatif
\[
\begin{tikzcd}[row sep=large, column sep=large]
 & X_0 \arrow[dl, "\alpha_1"'] \arrow[dr, "\alpha_2"] \arrow[dd, dotted, "i_0", ""'{name=J0, pos=0.6}] & \\
X_1 \arrow[dr, "i_1"', ""{name=J1, pos=0.4}] & & X_2 \arrow[dl, "i_2", ""'{name=J2, pos=0.4}] \\
 & \tX, &
\arrow[from=J0, to=J1, dotted, bend right, "\tilde{\lambda}_1"', shorten <=2pt, shorten >=2pt]
\arrow[from=J0, to=J2, dotted, bend left, "\tilde{\lambda}_2", shorten <=2pt, shorten >=2pt]
\end{tikzcd}
\]
% foot: Version 7 décembre 2002
\printedpage{92}
% head: A. Grothendieck : Dérivateurs, Ch. XII, Édition des Universités de Montpellier II et Paris VII
les flèches
\[
X_1 \longrightarrow \tX/1, \qquad X_2 \longrightarrow \tX/2
\]
étant dans $W$ (et même $W$-coasphériques - elles ont même un adjoint à gauche).

Donc on trouve, comme précédemment
\[
\left\{
\begin{aligned}
\alpha_1 \in W &\quad\iff\quad i_2 \in W \\
\alpha_2 \in W &\quad\iff\quad i_1 \in W,
\end{aligned}
\right.
\]
et de même on peut formuler un ‘lemme des cinq’ comme le corollaire page 125 \texttt{[plutôt 126]}, mais avec $\tX$, $\tX'$ au lieu de $\sX$, $\sX'$.

De plus, $\tX$ est la solution d'un problème universel dans Cat, en termes de données
\[
\begin{tikzcd}[row sep=large, column sep=large]
 & X_0 \arrow[dl, "\alpha_1"'] \arrow[dr, "\alpha_2"] \arrow[dd, "f_0", ""'{name=G0, pos=0.6}] & \\
X_1 \arrow[dr, "f_1"', ""{name=G1, pos=0.4}] & & X_2 \arrow[dl, "f_2", ""'{name=G2, pos=0.4}] \\
 & Y, &
\arrow[from=G0, to=G1, bend right, "\tilde{\lambda}_1"', shorten <=2pt, shorten >=2pt]
\arrow[from=G0, to=G2, bend left, "\tilde{\lambda}_2", shorten <=2pt, shorten >=2pt]
\end{tikzcd}
\]

\msp{page 136}
la seule différence avec celui formulé à la page 128 étant la direction opposée des flèches $\lambda_1$, $\lambda_2$. Pour tout carré commutatif
\[
\begin{tikzcd}
 & X_0 \arrow[dl, "\alpha_1"'] \arrow[dr, "\alpha_2"] & \\
X_1 \arrow[dr, "f_1"'] & & X_2 \arrow[dl, "f_2"] \\
 & Y, &
\end{tikzcd}
\tag*{$(*)$}
\]
on trouve encore une flèche canonique
\[
\tilde{\varphi} : \tX \longrightarrow Y.
\]

Il serait tout aussi raisonnable de dire que le carré $(*)$ est $W$-homotopiquement cocartésien, lorsque $\tilde{\varphi}$ est \texttt{[dans]} $W$, que d'adopter la définition 3, page 131 - on trouverait exactement les mêmes corollaires. Les deux définitions seraient elles équivalentes?? % [sic] seraient elles -> seraient-elles, hyphen of the inverted subject missing
% foot: Version 7 décembre 2002

\end{document}
